% Section 3: Methodology (1,000 words target)

\subsection{Data Preparation}

\subsubsection{Census Tracts (Baseline)}

We use census tract shapefiles and demographics from the 2020 Census Redistricting Data (PL 94-171 Summary File), yielding approximately 85,000 tracts nationally. Adjacency graphs were constructed using Rook contiguity (shared edges, excluding corner-only contact).

\subsubsection{Block Groups (Finer Resolution)}

Block group shapefiles and demographics were obtained from the Census Bureau TIGER/Line files. With approximately 240,000 block groups nationally ($\sim$3$\times$ more units than tracts), this resolution provides intermediate granularity. Adjacency graphs were constructed identically to the tract-level approach.

\subsubsection{Census Blocks (Finest Resolution)}

Census blocks represent the finest available resolution with approximately 11 million blocks nationally ($\sim$130$\times$ more units than tracts). Block-level shapefiles were obtained from TIGER/Line files, with demographics aggregated from the PL 94-171 file. This resolution poses significant computational challenges due to graph size (estimated 35 million edges nationally).

% TODO: Add memory and disk space requirements after implementation

\subsection{Algorithm Configuration}

To isolate resolution effects, we maintain identical parameters across all spatial scales:

\begin{itemize}
    \item \textbf{Population balance}: $\pm 0.5\%$ deviation from ideal district size
    \item \textbf{Edge-weighting}: $\alpha = 5.0$ weight multiplier for minority-minority edges
    \item \textbf{Minority threshold}: $\tau = 0.40$ (units with $\geq 40\%$ minority population)
    \item \textbf{Partitioning method}: METIS 5.1.0 recursive bisection with multilevel coarsening
\end{itemize}

This ensures that any observed differences arise from spatial resolution rather than parameter variation.

\subsection{Experimental Design}

We conduct redistricting for all 50 states at three resolutions (150 total runs) using 2020 Census data. To validate computational feasibility, we first test a 10-state subset:

\begin{itemize}
    \item \textbf{High minority}: Texas, California, Florida, Georgia, New York
    \item \textbf{Low minority}: Vermont, Maine, West Virginia, New Hampshire, Idaho
\end{itemize}

If subset results confirm METIS scalability at block-level resolution, we proceed to the full 50-state analysis.

\subsection{Metrics}

For each state and resolution, we compute:

\begin{itemize}
    \item \textbf{Majority-minority (MM) district count}: Districts where minority population $> 50\%$
    \item \textbf{Polsby-Popper compactness}: $\frac{4\pi \cdot \text{Area}}{\text{Perimeter}^2}$ (higher is more compact)
    \item \textbf{Population deviation}: Verify $\pm 0.5\%$ constraint maintained
    \item \textbf{Runtime}: Computational performance and scalability analysis
\end{itemize}

% TODO: Add statistical testing methodology (paired t-tests, regression models)
